\section{A quadratic physical cap at the critical gradient}
\label{sec:v98-cap}
The finite-count spectral error in \eqref{eq:v94-critical-finite}
has only a square-root gain in the gradient cutoff. Integrating that
error against the first inverse gradient is not legitimate. We obtain
a second, quadratic bound from the two-dimensional physical map. Its
constant is an explicitly weighted multiplicity, retained throughout
rather than replaced by an unestimated word count.

Fix $0<b_0<1/8$ within the range of \eqref{eq:v94-critical-finite}.
On the regular exact event put $G(x)=R(-p_0(x),p_m(x))$ and $z=|G|$.
The retained endpoint Hessian satisfies
$\lambda_\circ I\le K=D_q^2F\le\Lambda_\circ I$ in endpoint angles
$q=(\alpha_0,\alpha_m)$. The optical normalization is
$\omega=|\partial_{p_0}\alpha_m|^{-1}$.
Define, with physical states counted once,
\begin{align}
 D_{m,R,n,k}(v)
  &=\sum_{x\in\Pi^{\rm reg}:G(x)=v}
       \frac{I_m(x)\omega(x)}{4\pi c_*\det K(x)},\label{eq:v98-gradient-density}\\
 \mathcal K_m
  &=1+\pi\sup_{R,n,k}\operatorname*{ess\,sup}_{|v|<b_0}
                                      D_{m,R,n,k}(v).
                                      \label{eq:v98-weighted-multiplicity}
\end{align}
These definitions involve neither the roof target nor an inverse
power of $z$. The factor $I_m$ is the full optical insertion from
\eqref{eq:v94-derivative}.

\begin{lemma}[Physical area formula and a finite weighted cap]
\label{lem:v98-quadratic-cap}
For every finite $m$, $\mathcal K_m<\infty$, uniformly over the
permitted radius and exact labels. For $0<b<b_0$ and every interval
$J$,
\begin{equation}\label{eq:v98-quadratic-cap}
        \int_J u_b(t)\,dt\le\mathcal K_m b^2.
\end{equation}
More precisely, if $N_m$ bounds the cardinalities of the regular
physical fibers of $G$ at count $m$, and $v_*=2\tau_-/R_+$, then
\begin{equation}\label{eq:v98-finiteness-bound}
 \mathcal K_m\le1+C N_m(1+v_*)^{-(m-1)}.
\end{equation}
Equation~\eqref{eq:v98-finiteness-bound} establishes finiteness only;
no collision-uniform estimate of its right side is assumed.
\end{lemma}
\begin{proof}
First variation gives $F_{\alpha_0}=-Rp_0$. Change of variables from
$(\alpha_0,p_0)$ to endpoint angles has density $\omega$. The
Jacobian of $q\mapsto G(q)$ is $\det K>0$. Thus the area formula,
applied on injective local pieces and summed without multiplicity of
charts, gives
\[
 G_\#(I_m\bar\nu)=D_{m,R,n,k}(v)\,dv.
\]
Equivalently $|\det D_{(\alpha_0,p_0)}G|=\det K/\omega$.
This identity uses the original collision normalization and counts
physical preimages, not representatives in overlapping charts.
Truncation of $I_m$, followed by monotone convergence, justifies the
formula before its boundedness has been established.

Here is a bound for each summand uniform through internal grazing.
For flight lengths $L_i\ge\tau_-$ the retained positive optical
recursion is
\[
 y_1=L_0,\quad f_0=1,\quad
 f_i=f_{i-1}+\frac{2y_i}{Rc_i},\quad
 y_{i+1}=y_i+L_i f_i,\quad
 \omega=\frac{Rc_0c_m}{y_m}.
\]
Consequently
\[
 y_m\ge\tau_-\prod_{i=1}^{m-1}(1+v_*/c_i).
\]
For a selected internal incidence,
\[
 \frac1{c_i}\prod_j(1+v_*/c_j)^{-1}
 \le v_*^{-1}(1+v_*)^{-(m-2)}.
\]
Since $\sum_iq_0^{d_i}\le2q_0/(1-q_0)$, it follows that
$I_m\omega\le C(1+v_*)^{-(m-1)}$. For one flight the empty
product gives the same conclusion after enlarging $C$.
Division by $\det K\ge\lambda_\circ^2$ bounds each summand.

At a fixed count, the finite-horizon restriction permits only finitely
many lifted scatterer itineraries. The retained finite definable
physical graph, expressed in finitely many circle charts, includes
the first-hit tests, radius parameter and section decisions. The
nonzero Jacobian makes every regular $G$-fiber discrete. Uniform
finiteness for definable families therefore gives a finite $N_m$
\cite{Coste}. Displacement and occupation take finitely many values
at this count. This is used only for fixed-count finiteness and proves
\eqref{eq:v98-finiteness-bound}. Finally discard the roof test on the
upper side, and integrate \eqref{eq:v98-gradient-density} over the
disk $|v|<b$. Its area is $\pi b^2$, proving
\eqref{eq:v98-quadratic-cap}. No change is made to the source being
estimated.
\end{proof}

\begin{lemma}[Integration of two different positive caps]
\label{lem:v98-two-caps}
For $E,H,b>0$,
\begin{align}
 &b^{-1}\min(Eb^{1/2},Hb^2)
       +\int_0^b v^{-2}\min(Ev^{1/2},Hv^2)\,dv\notag\\
 &\hspace{25mm}\le3\min\{Hb,E^{2/3}H^{1/3}\}.
                                      \label{eq:v98-two-caps}
\end{align}
\end{lemma}
\begin{proof}
The crossing is $v_*=(E/H)^{2/3}$. For $b\le v_*$ the left side
is $2Hb$. For $b>v_*$ it is
$3E^{2/3}H^{1/3}-Eb^{-1/2}$. These formulas prove the claim.
The minimum is a cap on the positive source estimate, not on a
signed spectral remainder considered by itself.
\end{proof}

Define the actual untruncated source
\[
 w_b(t)\,dt=F_\#(I_mz^{-1}\one_{\{0<z<b\}}\bar\nu).
\]
It has a finite, absolutely continuous roof measure at every count.

\begin{theorem}[First inverse gradient with a complete finite-count error]
\label{thm:v98-uncut-insertion}
For one fixed auxiliary band $B_*$, all $m\ge2$, $0<b<b_0$ and
intervals $J$ of length $L$,
\begin{align}
 m^2\int_J w_b
 \le{}&C b^{\alpha-1}(L+B_*^{-1})\notag\\
 &+C\min\left\{m^2\mathcal K_m b,\ 
 C_{B_*}^{2/3}(1+B_*L)^{2/3}m^2
              (\mathcal K_m\theta_{B_*}^{2m})^{1/3}\right\}.
                                      \label{eq:v98-uncut-insertion}
\end{align}
The integration goes to gradient zero at finite count. In particular
there is no error proportional to a negative power of a lower cutoff.
The weighted multiplicity $\mathcal K_m$ is not asserted bounded or
subexponential in $m$.
\end{theorem}
\begin{proof}
Let $M(v)=m^2\int_Ju_v$. The two positive bounds give
\[
 M(v)\le C(L+B_*^{-1})v^\alpha
   +\min\{Ev^{1/2},Hv^2\},
 \quad E=C_{B_*}(1+B_*L)m^2\theta_{B_*}^m,
 \quad H=m^2\mathcal K_m.
\]
Indeed a nonnegative quantity bounded by both $A+D$ and $V$ is
bounded by $A+\min(D,V)$. Positive layer cake gives exactly
\[
 m^2\int_Jw_b=b^{-1}M(b)+\int_0^b v^{-2}M(v)\,dv.
\]
Apply it first above a positive lower endpoint, and then use monotone
convergence. The leading term is integrable because $\alpha>1$;
its multiplier is $1+(\alpha-1)^{-1}$. Lemma~\ref{lem:v98-two-caps}
integrates the full error. This proves the inequality and finiteness.
Absolute continuity follows from coarea on compact regular roof
charts and exhaustion; the finite critical set has zero source area.
\end{proof}

We isolate the additional estimate which would make this error
negligible. For an unbounded sequence of previously fixed auxiliary
bands, consider the condition
\begin{equation}\label{eq:v98-spectral-capacity}
       m^6\mathcal K_m\theta_{B_*}^{2m}\longrightarrow0.
                         \tag{SC}
\end{equation}
Subexponential growth of $\mathcal K_m$ is sufficient, but is not
proved here. Under \eqref{eq:v98-spectral-capacity}, first taking the
collision limit and only then the indicated band sequence gives
\begin{equation}\label{eq:v98-uncut-limit}
 \limsup_m\sup_{R,n,k}\sup_{|J|=L}m^2\int_Jw_b
                          \le C Lb^{1/128}.
\end{equation}
Adding $z\ge b_0$ by \eqref{eq:v94-cutoff} gives the bound $CL$
for $v^0=F_\#(I_mz^{-1}\bar\nu)/dt$ at each fixed $L$.
These are implications of the explicit condition, not unconditional
zero-cutoff conclusions. Geometric-pressure results for stable
Jacobian partition sums \cite{BDPressure} concern a different
weighted sum; no identification with \eqref{eq:v98-weighted-multiplicity}
or uniform parameter bound is imported from them.
